La nau \textit{BepiColombo}, de 2\,700~kg de massa, sobrevola Mercuri per estudiar el planeta.

\begin{apartat}{1,25}
Completeu la taula següent calculant la distància respecte al centre de Mercuri per als dos punts indicats.

\begin{center}
\begin{tabular}{|l|p{0.45\linewidth}|c|}
\cline{2-3}
\multicolumn{1}{c|}{} & \multicolumn{1}{c|}{\textit{Distància respecte al centre de Mercuri (m)}} & \textit{Energia potencial (J)} \\
\hline
\rule{0pt}{3.5ex}\textit{Punt 1} & & $-1,89\times 10^{10}$ \\[1.5ex]
\hline
\rule{0pt}{3.5ex}\textit{Punt 2} & & $-1,35\times 10^{10}$ \\[1.5ex]
\hline
\end{tabular}
\end{center}

En el punt més proper d'una altra aproximació a Mercuri, la intensitat del camp gravitatori és de $3,23\ \mathrm{m\,s^{-2}}$. Calculeu a quants quilòmetres de la superfície de Mercuri ha arribat la \textit{BepiColombo}.
\end{apartat}

\begin{apartat}{1,25}
Una part de la nau \textit{BepiColombo}, anomenada \textit{Mio}, està previst que se separi del conjunt i orbiti al voltant de Mercuri seguint una òrbita el·líptica, en què el punt més proper a Mercuri (punt $P$) es troba a 590~km de la superfície del planeta i el punt més llunyà a Mercuri (punt $A$) es troba a una distància d'11\,600~km sobre la superfície del planeta. La velocitat orbital de \textit{Mio} serà de $2,64\times 10^{3}\ \mathrm{m\,s^{-1}}$ en el punt $A$. Utilitzant el principi de conservació del moment angular, calculeu quin serà el mòdul de la velocitat de \textit{Mio} en el punt $P$. Justifiqueu si l'acceleració normal variarà al llarg de l'òrbita. En cas afirmatiu, indiqueu en quin punt tindrà el valor màxim.
\end{apartat}

\dades{%
  $G = 6,67\times 10^{-11}\ \mathrm{N\,m^2\,kg^{-2}}$.\\
  Massa de Mercuri: $M_\mathrm{M} = 3,285\times 10^{23}\ \mathrm{kg}$.\\
  Radi de Mercuri: $R_\mathrm{M} = 2,44\times 10^{3}\ \mathrm{km}$.}
